\chapter{Experiment}
%
\label{chap:Experiment}
%
This chapter describes the end-to-end experimental setup for evaluating LLMs
on the benchmark dataset introduced in Section~\ref{sec:benchmark}.  The
chapter covers the implementation of the extraction pipeline
(Section~\ref{sec:pipeline-impl}), the two evaluation protocols
(Section~\ref{sec:eval-protocols}), and the metrics used to compare models
(Section~\ref{sec:metrics}).

%----------------------------------------------------------------------
\section{Pipeline Implementation}
\label{sec:pipeline-impl}
%
\subsection{extraction.py}
%
The core pipeline is implemented in Python~3 as a single script~(\texttt{extraction.py}).
At a high level, the script:

\begin{enumerate}
    \item Reads the benchmark CSV~(250 rows, one row per ticket).
    \item For each row, locates the corresponding \texttt{.msg.md} file under
          the ticket folder by matching the ticket identifier in the
          \texttt{parent\_name} and \texttt{ticket\_number} columns.
    \item Concatenates all \texttt{.msg.md} files for that ticket into a
          single context string~(cf.\ Section~\ref{sec:data-prep}).
    \item Sends two sequential HTTP calls to the Ollama API
          (extraction call, then tagging call).
    \item Parses the JSON responses and writes the extracted fields back to the
          output CSV.
    \item Saves a checkpoint every~10 rows~(\texttt{SAVE\_EVERY~=~10}) to
          prevent data loss if the process is interrupted.
\end{enumerate}

A one-second sleep is inserted between rows to reduce pressure on the GPU
server~\texttt{ve-ds-nns-01:11434}.

\subsection{Content Loading}
%
Each ticket folder may contain multiple \texttt{.msg.md} files corresponding
to the email thread.  The function \texttt{load\_ticket\_content\_from\_md()}
iterates the folder and builds the context by appending each file's content
under a~\verb|[FILE: <name>]| header.  This preserves the chronological
sequence of the email exchange and ensures the model receives the complete
dialogue rather than just metadata.

\subsection{Ollama API Call}
%
The function \texttt{call\_ollama()} wraps the Ollama~\texttt{/api/generate}
endpoint with the following fixed parameters:

\begin{itemize}
    \item \texttt{model}: one of the ten configurations listed in
          Table~\ref{tab:models}
    \item \texttt{temperature}: $0.2$
    \item \texttt{num\_ctx}: $131{,}072$ (128k context window)
    \item \texttt{num\_predict}: $4{,}096$ (maximum output tokens)
    \item \texttt{stream}: \texttt{false}
\end{itemize}

A timeout of~180~seconds is set on the HTTP request to avoid indefinitely
blocking on slow models.

%----------------------------------------------------------------------
\section{Evaluation Protocols}
\label{sec:eval-protocols}
%
\subsection{Protocol 1 -- Blind Tag Evaluation}
\label{subsec:blind-eval}
%
In the blind tag evaluation, the model is given the raw ticket text and the
approved tag taxonomy~(117 tags) and is asked to select the most applicable
tags without any prior knowledge of the human annotations.  The model output
is compared to the ground-truth tags using token-level F1 (see
Section~\ref{sec:metrics}).

The evaluation is run separately for the \textbf{GS}, \textbf{BL}, and
\textbf{LB} categories, yielding three F1 scores per model.  A combined
macro-average F1 across all three categories is computed as the primary
ranking metric.

\subsection{Protocol 2 -- LLM-as-Judge Evaluation}
\label{subsec:judge-eval}
%
Because F1 alone may undervalue semantically correct but lexically different
tags~\cite{zheng2024judging}, a second evaluation track uses an LLM judge.
The judge model~(GPT-4o) receives:

\begin{enumerate}
    \item The original ticket text.
    \item The human ground-truth tags and Problem/Solution annotation.
    \item The model-generated tags, Problem, Solution, and reasoning.
\end{enumerate}

It then rates each model output on two dimensions on a 1--5 Likert scale:

\begin{description}
    \item[\textbf{Tag Usefulness~(TU)}:] Do the generated tags meaningfully
          describe the ticket content, even if they differ from the human
          labels?
    \item[\textbf{Reasoning Coherence~(RC)}:] Does the reasoning field
          demonstrate a logical, accurate understanding of the ticket?
\end{description}

Judge ratings are averaged over all 250 tickets per model.  To guard against
positional bias, the order in which model outputs are presented to the judge
is randomised per ticket.

%----------------------------------------------------------------------
\section{Evaluation Metrics}
\label{sec:metrics}
%
\subsection{Token-level F1}
%
For the blind tag evaluation, each tag is tokenised at the word level.  The
predicted tag set and ground-truth tag set are compared using standard
precision~$P$, recall~$R$, and F1:

\begin{equation}
    F_1 = \frac{2 \cdot P \cdot R}{P + R}
\end{equation}

where $P$ is the fraction of predicted tags found in the ground truth and $R$
is the fraction of ground-truth tags found in the predictions.  Macro-averaged
F1 across the three categories is used for the final ranking.

\subsection{Over- and Under-tagging Rates}
%
Beyond F1, two additional indicators are reported:

\begin{itemize}
    \item \textbf{Over-tagging rate}: Fraction of predicted tags that do not
          appear in the ground truth~(= $1-P$).
    \item \textbf{Under-tagging rate}: Fraction of ground-truth tags missed
          by the model~(= $1-R$).
\end{itemize}

These indicators reveal whether a model errs on the side of being too verbose
or too conservative.

\subsection{LLM-as-Judge Scores}
%
Mean Tag Usefulness~($\overline{TU}$) and mean Reasoning Coherence~($\overline{RC}$)
are reported per model on a 1--5 scale.  A composite judge score is defined
as the arithmetic mean:

\begin{equation}
    J = \frac{\overline{TU} + \overline{RC}}{2}
\end{equation}

Models are ranked separately by F1 and by $J$ to assess whether the two
protocols agree on the best-performing system.

\section{Validation Approach}
%
\subsection{Data Quality Checks}
%
Several validation steps were implemented:
\begin{enumerate}
    \item Verification of successful document conversion
    \item Validation of metadata extraction accuracy
    \item Consistency checks in the database
    \item Spot-checking of language detection results
\end{enumerate}

\subsection{Performance Monitoring}
%
System performance metrics were monitored including:
\begin{itemize}
    \item Processing time per document
    \item Memory usage during batch processing
    \item Error rates and failure modes
    \item Storage efficiency of converted documents
\end{itemize}

\section{Experimental Phases}
%
The experiment was conducted in multiple phases:

\begin{description}
    \item[Phase 1: Initial Processing] Convert all documents and extract basic metadata
    \item[Phase 2: Detailed Analysis] Perform language detection and content analysis
    \item[Phase 3: Aggregation] Aggregate results by document type and time period
    \item[Phase 4: Visualization] Generate comprehensive analysis charts and dashboards
\end{description}

\section{Tools and Environment}
%
\subsection{Development Environment}
%
Experiments were conducted in:
\begin{itemize}
    \item Jupyter Notebooks for interactive analysis
    \item Python virtual environment with controlled dependencies
    \item Windows environment with appropriate resource allocation
\end{itemize}

\subsection{Database Tools}
%
SQLite database operations were performed using:
\begin{itemize}
    \item Direct SQL queries for data retrieval
    \item Python sqlite3 library for programmatic access
    \item Database inspection tools for validation
\end{itemize}

\section{Reproducibility}
%
To ensure reproducibility:
\begin{itemize}
    \item All processing code is documented and version-controlled
    \item Processing parameters are explicitly documented
    \item Seed values for random processes are fixed
    \item Input dataset characteristics are recorded
    \item Output artifacts (databases, visualizations) are archived
\end{itemize}
